\documentclass[10pt,a4paper]{amsart}
\usepackage[margin=19mm]{geometry}
\usepackage{amsmath,amssymb,mathtools}
\usepackage[scr=rsfs,cal=euler]{mathalfa}
\usepackage{xfrac}
\usepackage[unicode,hidelinks,bookmarks=false]{hyperref}
\newcommand{\ie}{i.e.\ }
\newcommand{\cf}{cf.\ }
\newcommand{\resp}{respectively\ }
\newcommand{\Subsection}{\subsection}
\providecommand{\nobreakdash}{\nobreak\hbox{-}}
\setlength{\emergencystretch}{2em}
\multlinegap0pt
\usepackage{amssymb}
\usepackage{enumerate}
\newtheorem{proposition}{Proposition}
\title{Homogenization of Viscous Sublinear Hamilton--Jacobi Equations with $u/\epsilon$-Dependence}
\author{Guyu Jin}
\date{Source: arXiv:2608.00438v1 (CC BY 4.0)}
\begin{document}
\maketitle

\noindent\emph{Source and licence.} Homogenization of Viscous Sublinear Hamilton--Jacobi Equations with $u/\epsilon$-Dependence. Version arXiv:2608.00438v1 (CC BY 4.0), distributed by arXiv under the Creative Commons Attribution 4.0 International licence, http://creativecommons.org/licenses/by/4.0/, as recorded in the arXiv licence metadata for this exact version and shown on the abstract page https://arxiv.org/abs/2608.00438v1. Full licence notice: \texttt{licenses/arxiv-2608.00438-CC-BY-4.0.txt}.\par\smallskip

\noindent\textit{Editorial scope.} Counted unit: the article's proposition proposition (source0.tex lines 3841--3869). The article's complete original proof of it is delivered with it (source0.tex lines 3871--4454, one \texttt{proof} environment of 584 lines). No mathematics is written, repaired or re-derived: the statement and the proof are the article's own lines, in the article's own notation, and no retained line is substituted or re-flowed.  No further source-local context is needed: every cross-reference of every retained line resolves inside the two delivered spans.  Nothing was omitted from any retained span, so the package carries no omission record.  No retained line contains a citation, so the wrapper carries no bibliography and no bibliography entry.  No retained line contains a figure environment, an image-inclusion command or drawing code, so no image is delivered and the package declares no asset dependency.  Editorial formatting only, in addition: an AMS \texttt{amsart} preamble that restates the article's own declarations and macros used by the retained lines, namely the environments \texttt{proposition} on separate counters, together with the standard packages those lines need.  The retained environments print in this document as Proposition 1 in order of appearance, whereas the article prints the same items with the numbering of its own full document.  The complete native source of the article as distributed in the arXiv e-print for this exact version is bundled unchanged as \texttt{sources/206542/source0.tex}.
\begin{proposition}[Comparison principle in the linear growth class]
Let \(T>0\). Assume that
\[
 A(x)=\sigma(x)\sigma(x)^T,
 \qquad
 \sigma\in W^{1,\infty}(\mathbb R^n;\mathbb R^{n\times m}).
\]
Assume that \(H=H(x,r,p)\) is continuous and satisfies the following Lipschitz condition: there exists \(L>0\) such that
\[
 |H(x,r,p)-H(y,s,q)|
 \le
 L\Bigl((1+|p|+|q|)|x-y|+|r-s|+|p-q|\Bigr)
\]
for all \(x,y\in\mathbb R^n\), \(r,s\in\mathbb R\), and \(p,q\in\mathbb R^n\). Let \(u\in C(\mathbb R^n\times[0,T])\) be a viscosity subsolution and \(v\in C(\mathbb R^n\times[0,T])\) be a viscosity supersolution of
\[
 w_t+H(x,w,Dw)=\operatorname{tr}(A(x)D^2w)
 \quad\text{in }\mathbb R^n\times(0,T).
\]
Assume that
\[
 u(x,0)\le v(x,0)
 \quad\text{for all }x\in\mathbb R^n,
\]
and that \(u,v\) have at most linear growth in \(x\), uniformly in \(t\in[0,T]\). Then
\[
 u(x,t)\le v(x,t)
 \quad\text{for all }(x,t)\in\mathbb R^n\times[0,T].
\]
\end{proposition}

\begin{proof}
Choose \(k>L\) and set
\[
 \widetilde u(x,t):=e^{-kt}u(x,t),
 \qquad
 \widetilde v(x,t):=e^{-kt}v(x,t).
\]
Then \(\widetilde u\) and \(\widetilde v\) are respectively a viscosity subsolution and supersolution of
\[
 w_t+G(t,x,w,Dw)=\operatorname{tr}(A(x)D^2w),
\]
where
\[
 G(t,x,r,p):=kr+e^{-kt}H(x,e^{kt}r,e^{kt}p).
\]
Moreover, if \(r\ge s\), then
\[
\begin{aligned}
 G(t,x,r,p)-G(t,x,s,p)
 &=
 k(r-s)+e^{-kt}
 \Bigl[
 H(x,e^{kt}r,e^{kt}p)-H(x,e^{kt}s,e^{kt}p)
 \Bigr] \\
 &\ge (k-L)(r-s).
\end{aligned}
\]
Set $ \lambda:=k-L>0.$

We prove \(\widetilde u\le \widetilde v\). Suppose by contradiction that $M:=\sup_{\mathbb R^n\times[0,T]}(\widetilde u-\widetilde v)>0.$

We first introduce a localization function. Since \(\widetilde u,\widetilde v\) have at most linear growth, there exists \(C_0>0\) such that
\[
 \widetilde u(x,t)-\widetilde v(y,s)\le C_0(1+|x|+|y|)
 \quad\text{for all }x,y\in\mathbb R^n,\ t,s\in[0,T].
\]
Choose a family \(\{\beta_R\}_{R>1}\subset C^2(\mathbb R^n)\) such that
\[
 \beta_R\ge0,\qquad
 \beta_R(x)=0\quad\text{for }|x|\le R,
\]
\[
 \liminf_{|x|\to\infty}\frac{\beta_R(x)}{|x|}\ge 4C_0,
\]
\[
 \|D\beta_R\|_{L^\infty}+\|D^2\beta_R\|_{L^\infty}\le C,
\]
and
\[
 \beta_R(x)\to0
 \quad\text{as }R\to\infty
 \quad\text{for every fixed }x\in\mathbb R^n.
\]
\textbf{Step 1:}

Let $ \zeta(z):=\sqrt{1+|z|^2}.$
We first prove the auxiliary estimate
\[
 \widetilde u(x,t)-\widetilde v(y,t)
 \le
 C(1+|x-y|)
 \quad\text{for all }x,y\in\mathbb R^n,\ t\in[0,T].
\]
Indeed, fix \(\Theta>0\), \(\mu>0\), and \(R>1\), and consider
\[
 \Phi_R(x,y,t)
 :=
 \widetilde u(x,t)-\widetilde v(y,t)
 -
 \Theta e^{\mu t}\zeta(x-y)
 -
 \beta_R(x)-\beta_R(y).
\]
Because of the growth of \(\beta_R\), \(\Phi_R\) attains its supremum on \(\mathbb R^n\times\mathbb R^n\times[0,T]\). We claim that for \(\Theta,\mu\) sufficiently large, independent of \(R\), $ \sup_{\mathbb R^n\times\mathbb R^n\times[0,T]}\Phi_R\le0.$
Assume otherwise. By the initial condition and by choosing \(\Theta\) larger if necessary, a positive maximum cannot occur at \(t=0\). To avoid the common time variable, for \(\nu>0\) consider
\[
 \Psi_{R,\nu}(x,y,t,s)
 :=
 \widetilde u(x,t)-\widetilde v(y,s)
 -
 \Theta e^{\mu t}\zeta(x-y)
 -
 \beta_R(x)-\beta_R(y)
 -
 \frac{|t-s|^2}{2\nu}.
\]
Let \((x_\nu,y_\nu,t_\nu,s_\nu)\) be a maximum point. As \(\nu\to0\), up to a subsequence, $(x_\nu,y_\nu,t_\nu,s_\nu)\to(\bar x,\bar y,\bar t,\bar t),$
where \((\bar x,\bar y,\bar t)\) is a positive maximum point of \(\Phi_R\).

We now derive the key estimate used in Step 1. Let
\[
 p_\nu:=\Theta e^{\mu t_\nu}D\zeta(x_\nu-y_\nu),
 \qquad
 Z_\nu:=\Theta e^{\mu t_\nu}D^2\zeta(x_\nu-y_\nu).
\]
By Ishii's lemma, there exist \(X_\nu,Y_\nu\in\mathbb S^n\) such that
\[
 \left(
 \mu\Theta e^{\mu t_\nu}\zeta(x_\nu-y_\nu)
 +
 \frac{t_\nu-s_\nu}{\nu},
 p_\nu+D\beta_R(x_\nu),
 X_\nu+D^2\beta_R(x_\nu)
 \right)
 \in
 \overline{J}^{2,+}\widetilde u(x_\nu,t_\nu),
\]
and
\[
 \left(
 \frac{t_\nu-s_\nu}{\nu},
 p_\nu-D\beta_R(y_\nu),
 Y_\nu-D^2\beta_R(y_\nu)
 \right)
 \in
 \overline{J}^{2,-}\widetilde v(y_\nu,s_\nu).
\]
Using the viscosity subsolution inequality for \(\widetilde u\), we get
\[
\begin{aligned}
 &\mu\Theta e^{\mu t_\nu}\zeta(x_\nu-y_\nu)
 +
 \frac{t_\nu-s_\nu}{\nu}
 +
 G\left(
 t_\nu,x_\nu,\widetilde u(x_\nu,t_\nu),
 p_\nu+D\beta_R(x_\nu)
 \right)
 \\
 &\le
 \operatorname{tr}
 \left(
 A(x_\nu)\bigl(X_\nu+D^2\beta_R(x_\nu)\bigr)
 \right).
\end{aligned}
\]
Similarly, using the viscosity supersolution inequality for \(\widetilde v\), we have
\[
\begin{aligned}
 &\frac{t_\nu-s_\nu}{\nu}
 +
 G\left(
 s_\nu,y_\nu,\widetilde v(y_\nu,s_\nu),
 p_\nu-D\beta_R(y_\nu)
 \right)
 \\
 &\ge
 \operatorname{tr}
 \left(
 A(y_\nu)\bigl(Y_\nu-D^2\beta_R(y_\nu)\bigr)
 \right).
\end{aligned}
\]
Subtracting the two inequalities yields
\[
\begin{aligned}
 &\mu\Theta e^{\mu t_\nu}\zeta(x_\nu-y_\nu)
 +
 G\left(
 t_\nu,x_\nu,\widetilde u(x_\nu,t_\nu),
 p_\nu+D\beta_R(x_\nu)
 \right)
 \\
 &\quad -
 G\left(
 s_\nu,y_\nu,\widetilde v(y_\nu,s_\nu),
 p_\nu-D\beta_R(y_\nu)
 \right)
 \\
 &\le
 \operatorname{tr}
 \left(
 A(x_\nu)\bigl(X_\nu+D^2\beta_R(x_\nu)\bigr)
 \right)
 -
 \operatorname{tr}
 \left(
 A(y_\nu)\bigl(Y_\nu-D^2\beta_R(y_\nu)\bigr)
 \right).
\end{aligned}
\]
We split the difference of the \(G\)-terms as
\[
\begin{aligned}
 &G\left(
 t_\nu,x_\nu,\widetilde u(x_\nu,t_\nu),
 p_\nu+D\beta_R(x_\nu)
 \right)
 -
 G\left(
 s_\nu,y_\nu,\widetilde v(y_\nu,s_\nu),
 p_\nu-D\beta_R(y_\nu)
 \right)
 \\
 &=
 \Bigl[
 G\left(
 t_\nu,x_\nu,\widetilde u(x_\nu,t_\nu),
 p_\nu+D\beta_R(x_\nu)
 \right)
 -
 G\left(
 t_\nu,x_\nu,\widetilde v(y_\nu,s_\nu),
 p_\nu+D\beta_R(x_\nu)
 \right)
 \Bigr]
 \\
 &\quad+
 \Bigl[
 G\left(
 t_\nu,x_\nu,\widetilde v(y_\nu,s_\nu),
 p_\nu+D\beta_R(x_\nu)
 \right)
 -
 G\left(
 s_\nu,y_\nu,\widetilde v(y_\nu,s_\nu),
 p_\nu-D\beta_R(y_\nu)
 \right)
 \Bigr].
\end{aligned}
\]
Since \(G\) is strictly increasing in the \(r\)-variable, we have
\[
\begin{aligned}
 &G\left(
 t_\nu,x_\nu,\widetilde u(x_\nu,t_\nu),
 p_\nu+D\beta_R(x_\nu)
 \right)
 -
 G\left(
 t_\nu,x_\nu,\widetilde v(y_\nu,s_\nu),
 p_\nu+D\beta_R(x_\nu)
 \right)
 \\
 &\ge
 \lambda\bigl(
 \widetilde u(x_\nu,t_\nu)-\widetilde v(y_\nu,s_\nu)
 \bigr),
\end{aligned}
\]
where \(\lambda=k-L>0\). Therefore
\[
\begin{aligned}
 \mu\Theta e^{\mu t_\nu}\zeta(x_\nu-y_\nu)
 &\le
 -\lambda\bigl(
 \widetilde u(x_\nu,t_\nu)-\widetilde v(y_\nu,s_\nu)
 \bigr)
 \\
 &\quad+
 E_\nu
 +
 \mathcal T_\nu,
\end{aligned}
\]
where
\[
\begin{aligned}
 E_\nu
 &:=
 \left|
 G\left(
 t_\nu,x_\nu,\widetilde v(y_\nu,s_\nu),
 p_\nu+D\beta_R(x_\nu)
 \right)
 -
 G\left(
 s_\nu,y_\nu,\widetilde v(y_\nu,s_\nu),
 p_\nu-D\beta_R(y_\nu)
 \right)
 \right|,
\end{aligned}
\]
and
\[
\begin{aligned}
 \mathcal T_\nu
 &:=
 \operatorname{tr}
 \left(
 A(x_\nu)\bigl(X_\nu+D^2\beta_R(x_\nu)\bigr)
 \right)
 -
 \operatorname{tr}
 \left(
 A(y_\nu)\bigl(Y_\nu-D^2\beta_R(y_\nu)\bigr)
 \right).
\end{aligned}
\]
We now estimate \(E_\nu\) and \(\mathcal T_\nu\). First, since $|D\zeta|\le1,
 $ $
 |p_\nu|\le \Theta e^{\mu T},$
and since $ \|D\beta_R\|_{L^\infty}\le C,$
the gradient variables in \(E_\nu\) are bounded by \(C(1+\Theta e^{\mu T})\). Moreover,
\[
 |D\beta_R(x_\nu)+D\beta_R(y_\nu)|\le C.
\]
Using the Lipschitz continuity of \(G\) in \(x\) and \(p\), we obtain
\[
 E_\nu
 \le
 C\Theta e^{\mu t_\nu}\zeta(x_\nu-y_\nu)+C+\omega_R(|t_\nu-s_\nu|),
\]
where \(\omega_R(r)\to0\) as \(r\to0\). The modulus \(\omega_R\) comes from the time-dependence of \(G\); for fixed \(R\), the variables are confined to a compact set by the localization \(\beta_R\).

Next, using \(A=\sigma\sigma^T\) and the matrix inequality in Ishii's lemma, we estimate the second-order term as follows:
\[
\begin{aligned}
 \mathcal T_\nu
 &\le
 C\Theta e^{\mu t_\nu}
 \|D^2\zeta(x_\nu-y_\nu)\|
 \|\sigma(x_\nu)-\sigma(y_\nu)\|^2
 \\
 &\quad+
 C\bigl(
 \|D^2\beta_R(x_\nu)\|
 +
 \|D^2\beta_R(y_\nu)\|
 \bigr).
\end{aligned}
\]
Since \(\sigma\) is Lipschitz and
\[
 \|D^2\zeta(z)\|\le C(1+|z|^2)^{-1/2},
\]
we get
\[
\begin{aligned}
 \mathcal T_\nu
 &\le
 C\Theta e^{\mu t_\nu}
 \frac{|x_\nu-y_\nu|^2}{\sqrt{1+|x_\nu-y_\nu|^2}}
 +
 C
 \\
 &\le
 C\Theta e^{\mu t_\nu}\zeta(x_\nu-y_\nu)+C.
\end{aligned}
\]
Combining the estimates for \(E_\nu\) and \(\mathcal T_\nu\), we arrive at
\[
\begin{aligned}
 \mu\Theta e^{\mu t_\nu}\zeta(x_\nu-y_\nu)
 &\le
 -\lambda\bigl(
 \widetilde u(x_\nu,t_\nu)-\widetilde v(y_\nu,s_\nu)
 \bigr)
 \\
 &\quad+
 C\Theta e^{\mu t_\nu}\zeta(x_\nu-y_\nu)
 +
 C
 +
 \omega_R(|t_\nu-s_\nu|).
\end{aligned}
\]
Finally, since \(t_\nu-s_\nu\to0\) as \(\nu\to0\), we may let \(\nu\to0\) and obtain
\[
\begin{aligned}
 \mu\Theta e^{\mu \bar t}\zeta(\bar x-\bar y)
 &\le
 -\lambda\bigl(
 \widetilde u(\bar x,\bar t)-\widetilde v(\bar y,\bar t)
 \bigr)
 \\
 &\quad+
 C\Theta e^{\mu \bar t}\zeta(\bar x-\bar y)
 +
 C.
\end{aligned}
\]
Moreover,
\[
 \mu\Theta e^{\mu\bar t}\zeta(\bar x-\bar y)
 \le
 C\Theta e^{\mu\bar t}\zeta(\bar x-\bar y)+C.
\]
Choosing \(\mu>2C\) and then \(\Theta\) sufficiently large yields a contradiction. Therefore $\Phi_R\le0.$
Letting \(R\to\infty\), we obtain
\[
 \widetilde u(x,t)-\widetilde v(y,t)
 \le
 \Theta e^{\mu T}\zeta(x-y)
 \le
 C(1+|x-y|).
\]

\textbf{Step 2:}

We now prove the comparison. Let \(\eta>0\), and assume that
\[
 \sup_{\mathbb R^n\times[0,T]}
 \bigl(\widetilde u-\widetilde v-\eta t\bigr)>0.
\]
For \(\alpha,\kappa,\nu>0\), consider
\[
\begin{aligned}
 M_{\alpha,\kappa,\nu}
 :=
 \sup_{x,y\in\mathbb R^n,\ t,s\in[0,T]}
 \Bigl\{
 \widetilde u(x,t)-\widetilde v(y,s)
 -
 \eta t
 -
 \alpha|x|^2
 -
 \frac{|x-y|^2}{2\kappa}
 -
 \frac{|t-s|^2}{2\nu}
 \Bigr\}.
\end{aligned}
\]
Let \((\bar x,\bar y,\bar t,\bar s)\) be a maximum point. The estimate just proved implies
\[
 |\bar x-\bar y|\le C\kappa,
\]
and the standard doubling argument gives
\[
 \frac{|\bar t-\bar s|^2}{\nu}\to0
 \quad\text{as }\nu\to0.
\]
Moreover, $ \alpha|\bar x|^2$
is bounded, hence
\[
 \alpha|\bar x|\to0
 \quad\text{as }\alpha\to0.
\]

For small parameters, the positivity of the above supremum and the initial condition imply
\[
 \bar t>0,
 \qquad
 \bar s>0.
\]
Set
\[
 p:=\frac{\bar x-\bar y}{\kappa}+2\alpha\bar x,
 \qquad
 q:=\frac{\bar x-\bar y}{\kappa}.
\]
By the parabolic Ishii lemma, there exist \(X,Y\in\mathbb S^n\) such that
\[
 \left(
 \eta+\frac{\bar t-\bar s}{\nu},
 p,
 X
 \right)
 \in
 \overline{J}^{2,+}\widetilde u(\bar x,\bar t),
\]
and
\[
 \left(
 \frac{\bar t-\bar s}{\nu},
 q,
 Y
 \right)
 \in
 \overline{J}^{2,-}\widetilde v(\bar y,\bar s).
\]
We estimate the second-order term. Since
\[
 A(x)=\sigma(x)\sigma(x)^T,
 \qquad
 \sigma\in W^{1,\infty}(\mathbb R^n;\mathbb R^{n\times m}),
\]
we write
\[
\begin{aligned}
&\operatorname{tr}(A(\bar x)X)-\operatorname{tr}(A(\bar y)Y)\\
&=
\sum_{j=1}^m
\left[
\left\langle X\sigma(\bar x)e_j,\sigma(\bar x)e_j\right\rangle
-
\left\langle Y\sigma(\bar y)e_j,\sigma(\bar y)e_j\right\rangle
\right].
\end{aligned}
\]
By the matrix inequality in Ishii's lemma,
\[
 \begin{pmatrix}
 X & 0\\
 0 & -Y
 \end{pmatrix}
 \le
 \frac{C}{\kappa}
 \begin{pmatrix}
 I & -I\\
 -I & I
 \end{pmatrix}
 +
 C\alpha
 \begin{pmatrix}
 I & 0\\
 0 & I
 \end{pmatrix}.
\]
Applying this inequality to the vector
\[
 \binom{\sigma(\bar x)e_j}{\sigma(\bar y)e_j}
\]
and summing over \(j=1,\dots,m\), we get
\[
\begin{aligned}
&\operatorname{tr}(A(\bar x)X)-\operatorname{tr}(A(\bar y)Y)\\
&\le
\frac{C}{\kappa}
\|\sigma(\bar x)-\sigma(\bar y)\|^2
+
C\alpha
\left(
\|\sigma(\bar x)\|^2+\|\sigma(\bar y)\|^2
\right).
\end{aligned}
\]
Since \(\sigma\) is bounded and Lipschitz continuous, this yields
\[
 \operatorname{tr}(A(\bar x)X)-\operatorname{tr}(A(\bar y)Y)
 \le
 C\frac{|\bar x-\bar y|^2}{\kappa}+C\alpha.
\]
Using the viscosity inequalities, we obtain
\[
\begin{aligned}
 \eta
 &+
 G(\bar t,\bar x,\widetilde u(\bar x,\bar t),p)
 -
 G(\bar s,\bar y,\widetilde v(\bar y,\bar s),q) \\
 &\le
 C\frac{|\bar x-\bar y|^2}{\kappa}+C\alpha.
\end{aligned}
\]
By the strict monotonicity of \(G\) in the \(r\)-variable,
\[
\begin{aligned}
 &G(\bar t,\bar x,\widetilde u(\bar x,\bar t),p)
 -
 G(\bar s,\bar y,\widetilde v(\bar y,\bar s),q)\\
 &\ge
 \lambda\bigl(\widetilde u(\bar x,\bar t)-\widetilde v(\bar y,\bar s)\bigr)
 -
 E_{\alpha,\kappa,\nu},
\end{aligned}
\]
where $ E_{\alpha,\kappa,\nu}
 :=
 \left|
 G(\bar t,\bar x,\widetilde v(\bar y,\bar s),p)
 -
 G(\bar s,\bar y,\widetilde v(\bar y,\bar s),q)
 \right|.$
Since
\[
 |\bar x-\bar y|\le C\kappa,
 \qquad
 |p-q|=2\alpha|\bar x|\to0,
\]
and \(p,q\) remain bounded, we have $E_{\alpha,\kappa,\nu}\to0$
as \(\nu\to0\), then \(\kappa\to0\), and finally \(\alpha\to0\). Hence
\[
 \eta
 \le
 C\frac{|\bar x-\bar y|^2}{\kappa}+C\alpha+E_{\alpha,\kappa,\nu}.
\]
Letting \(\nu\to0\), then \(\kappa\to0\), and then \(\alpha\to0\), we obtain $\eta\le0,$
which is impossible. Therefore, for every \(\eta>0\),
\[
 \widetilde u(x,t)-\widetilde v(x,t)-\eta t\le0
 \quad\text{in }\mathbb R^n\times[0,T].
\]
Letting \(\eta\to0\), we obtain
\[
 \widetilde u\le \widetilde v.
\]
Since \(e^{kt}>0\), this gives
\[
 u\le v
 \quad\text{in }\mathbb R^n\times[0,T].
\]
\end{proof}

\end{document}
